% Generated draft: requires mathematical and source review.
\documentclass[11pt]{article}
\usepackage[margin=25mm]{geometry}
\usepackage{fontspec}
\setmainfont{Latin Modern Roman}
\usepackage{amsmath,amssymb,mathtools,unicode-math}
\setmathfont{Latin Modern Math}
\usepackage{enumitem,needspace,xurl,hyperref}
\setlength{\parindent}{0pt}
\setlength{\parskip}{5pt}
\setlength{\emergencystretch}{3em}
\newcommand{\N}{\mathbb{N}}
\newcommand{\Z}{\mathbb{Z}}
\newcommand{\Q}{\mathbb{Q}}
\newcommand{\R}{\mathbb{R}}
\newcommand{\C}{\mathbb{C}}
\newcommand{\F}{\mathbb{F}}
\newcommand{\A}{\mathbb{A}}
\newcommand{\PP}{\mathbb{P}}
\newcommand{\E}{\mathbb{E}}
\newcommand{\eps}{\varepsilon}
\newcommand{\dd}{\,\mathrm{d}}
\begin{document}
\section*{UnsolvedMath ID: 800377}
\textbf{Source problem code:} catalogue entry
\subsection{Definitions and mathematical statement}
The supplied excerpt does not identify the matrix whose submatrices occur in the conjecture. Consequently, only the following statement schema can be recorded; it is not a complete mathematical formulation.

Write \(A=(a_{p,q})\) for the unspecified matrix. For \(i,j,r\in\mathbb{N}\), define, whenever the indicated entries exist,
\[
D_{j,r}(i)=\det\bigl(a_{i+s,j+t}\bigr)_{0\le s,t\le r}.
\]
Thus \(D_{j,r}(i)\) is the determinant of the \((r+1)\times(r+1)\) submatrix with consecutive rows \(i,i+1,\ldots,i+r\) and consecutive columns \(j,j+1,\ldots,j+r\).

The assertion attributed to Brenti, Conjecture 15 (2024), is that, for every fixed \(j,r\in\mathbb{N}\), the function \(i\mapsto D_{j,r}(i)\) is monotonically increasing on \(\mathbb{N}\). If increasing means nondecreasing, the asserted inequalities are
\[
D_{j,r}(i+1)\ge D_{j,r}(i)
\qquad\text{for all }i,j,r\in\mathbb{N}.
\]
If strict increase is intended, \(\ge\) must instead be replaced by \(>\).

The matrix, its entries and their domain, any additional parameters or hypotheses, and the ordering in which determinants are compared are absent from the supplied record. The convention for whether \(0\in\mathbb{N}\) is also unspecified; it affects both the allowed starting indices and whether \(r=0\) is included. These omissions must be resolved before the conjecture has a precise mathematical target.
\subsection{Short English statement}
For every fixed column starting index \(j\) and block size \(r+1\), the conjecture asserts that the determinant of a consecutive square block increases as its row starting index \(i\) increases. The supplied excerpt omits the underlying matrix and does not clarify whether increase is strict, so the statement remains incomplete.
\subsection{Sources}
\begin{itemize}
\item Supplied local UnsolvedMath registry record: permanent ID \(800377\), exact source problem code \texttt{catalogue entry}, title ``Coxeter Groups and Unimodality --- Conjecture 15.'' The record attributes the conjecture to Brenti (2024) and labels it open; that status is supplied curation, not an independently verified mathematical assessment.
\item Brenti, \emph{Some Open Problems on Coxeter Groups and Unimodality} (2024), Conjecture 15. Supplied source URL: \url{https://arxiv.org/abs/2410.09897}. The supplied background records access on 2026-10-16 and extraction by \texttt{arxiv\_latex}. These are supplied provenance details; no independent source visit or verification is claimed.
\end{itemize}
\end{document}
